%------------------------------------------------------------------------------
\subsection{Part I: The Coordinate System}
\label{sec:scaling}
%------------------------------------------------------------------------------

\subsubsection{Log-Frequency Scaling is Formally Demonstrated}
\label{sec:logscaling}

The claim that brain oscillation bands are logarithmically spaced has been repeated in the literature for over two decades \parencite{penttonen2003natural,buzsaki2004neuronal}, yet no study has submitted it to formal statistical test. We do so here using the simplest possible design: fit polynomial density models to the peak frequency distribution in both linear-Hz and log-Hz coordinates, then compare model quality with the Bayesian Information Criterion.

Across polynomial degrees 3 through 5, log-frequency description is decisively superior (Table~\ref{tab:polynomial_bic}). A degree-3 polynomial in log-Hz achieves BIC $= -1065.9$, better than a degree-5 polynomial in linear-Hz (BIC $= -1010.9$). The $\Delta$BIC advantage of log over linear ranges from 46.2 at degree 3 to 57.9 at degree 5. By the standard Kass--Raftery scale, $\Delta$BIC $> 10$ constitutes ``very strong'' evidence against the weaker model; our differences exceed 46. The peak distribution of 4,572,489 oscillatory peaks is not merely somewhat simpler in log-frequency space---it is fundamentally more parsimonious. Linear band spacing is not a viable description of the data.

\begin{table}[h]
\centering
\caption{Polynomial BIC comparison: log-Hz versus linear-Hz description of peak density.
Lower BIC indicates better fit penalized for model complexity.
$\Delta$BIC $>$ 10 is ``very strong'' evidence (Kass \& Raftery, 1995).}
\label{tab:polynomial_bic}
\begin{tabular}{r r r l r}
\toprule
Degree & Linear-Hz BIC & Log-Hz BIC & Winner & $\Delta$BIC \\
\midrule
3 & $-1019.8$ & $-1065.9$ & log & \textbf{46.2} \\
4 & $-1015.2$ & $-1072.5$ & log & \textbf{57.3} \\
5 & $-1010.9$ & $-1068.8$ & log & \textbf{57.9} \\
\bottomrule
\end{tabular}
\end{table}

\subsubsection{Band Boundaries are Genuine Spectral Features}
\label{sec:boundaries_genuine}

A standing concern in EEG band research is that the ``gaps'' between bands might be artefacts---either of the $1/f$ aperiodic background (which creates a falling density gradient that could be mistaken for troughs) or of researcher conventions that have influenced which frequencies are measured and reported. We test this with a Poisson surrogate null.

For each of 200 surrogate datasets, peaks were drawn from a Poisson process matched to the smooth density envelope (Gaussian smoothing $\sigma = 40$ bins on a 1000-bin log-frequency histogram), which preserves the overall $1/f$ decay while destroying all band structure. The deepest trough in any surrogate never fell below a depth ratio of 0.99---surrogates are uniformly flat at the $\pm 0.2\%$ level. In the real data, five troughs are observed with depth ratios of 0.296, 0.913, 0.383, 0.884, and 0.678 at 5.1, 7.8, 13.4, 25.3, and 35.0 Hz respectively (Table~\ref{tab:aperiodic_null}). Each trough is individually significant at $p < 0.0001$ against the surrogate distribution. The deepest trough (5.1 Hz, depth 0.296) represents a 70.4\% depletion relative to the smooth envelope---a spectral void that no aperiodic gradient can produce.

\begin{table}[h]
\centering
\caption{Aperiodic-only null test. Trough ``depth'' is the ratio of observed peak count to smooth envelope prediction; a ratio of 1.0 means no depletion. Surrogate distributions (200 Poisson samples from the smooth envelope) never produce ratios below 0.99, confirming that all five troughs are genuine band boundaries rather than $1/f$ artefacts.}
\label{tab:aperiodic_null}
\begin{tabular}{r r r r l}
\toprule
Trough (Hz) & Real depth & Surrogate mean $\pm$ SD & Depletion & $p$ \\
\midrule
5.1  & 0.296 & $1.004 \pm 0.004$ & 70.4\% & $< 0.0001$ \\
7.8  & 0.913 & $0.998 \pm 0.003$ & 8.7\%  & $< 0.0001$ \\
13.4 & 0.383 & $1.003 \pm 0.003$ & 61.7\% & $< 0.0001$ \\
25.3 & 0.884 & $1.000 \pm 0.002$ & 11.6\% & $< 0.0001$ \\
35.0 & 0.678 & $1.002 \pm 0.003$ & 32.2\% & $< 0.0001$ \\
\bottomrule
\end{tabular}
\end{table}

Additionally, the real trough ratio (geometric mean 1.621) lies 7.3 standard deviations below the surrogate distribution (mean 2.508, SD 0.122; $z = -7.29$). Aperiodic surrogates produce widely-spaced shallow undulations; the real data has tightly-spaced deep troughs. Band boundaries are genuine spectral features of oscillatory organisation, not analytical artefacts.

\begin{figure}[htbp]
\centering
\includegraphics[width=\textwidth]{fig1_peak_density.png}
\caption{\textbf{Peak frequency distribution and model comparison.} (A)~Gaussian KDE of 4,572,489 FOOOF-detected peaks from $N = 2{,}097$ subjects (9 datasets) in log-Hz space. Gold dashed lines mark $\phisym$-lattice boundaries ($\fzero \times \phisym^n$); red triangles mark empirical density troughs. The geometric mean ratio of consecutive troughs is 1.6175, indistinguishable from $\phisym = 1.618$. (B)~Bayesian Information Criterion (BIC) comparison of geometric series models fit to trough positions (least-squares in log-frequency space, 1 free parameter per fixed-ratio model, 2 for the free-ratio model). $\phisym$ achieves the lowest BIC ($-24.2$); linear spacing is excluded by $\Delta$BIC $> 15$.}
\label{fig:peak_density}
\end{figure}

\begin{figure}[htbp]
\centering
\includegraphics[width=\textwidth]{fig2_aperiodic_null.png}
\caption{\textbf{Aperiodic-only null test.} (A)~Real peak density (blue, $N = 4{,}572{,}489$ peaks) versus the smooth aperiodic envelope (red) obtained by heavy Gaussian smoothing ($\sigma = 40$ bins on 1,000-bin log-frequency histogram). Green dashed lines mark the five empirical troughs. (B)~Distribution of the deepest trough depth across 200 Poisson surrogates drawn from the aperiodic envelope. The real deepest trough (depth ratio $= 0.296$, red line) lies far outside the surrogate distribution (surrogate mean $= 1.004 \pm 0.004$ SD; permutation test, $p < 0.0001$), confirming that density troughs reflect genuine band structure, not $1/f$ artefacts.}
\label{fig:aperiodic_null}
\end{figure}

\subsubsection{Boundaries are Stable Across Smoothing Parameters}
\label{sec:stability}

The number and location of detectable troughs in a kernel density estimate depends on the smoothing bandwidth---over-smoothed estimates merge nearby troughs, while under-smoothed estimates fragment genuine troughs into noise. A robust boundary must survive across a wide range of smoothing levels.

We computed trough locations across 30 smoothing bandwidths ($\sigma = 2$--30 bins on the 1000-bin log-frequency histogram). Three troughs survive the full range of smoothing levels: $5.18 \pm 0.17$ Hz (30/30 smoothings), $13.44 \pm 0.05$ Hz (30/30), and $25.29 \pm 0.40$ Hz (24/30). These are the most bandwidth-robust troughs---features that no choice of smoothing can eliminate. Additional troughs at approximately 7.6 Hz and 35 Hz emerge at narrow smoothing bandwidths ($\sigma < 8$ bins) but merge into adjacent features at wider smoothings. All five troughs are confirmed genuine by the aperiodic null (Section~\ref{sec:boundaries_genuine}).

Note that bandwidth stability and trough depth measure different properties. The $\sim$35 Hz trough is deeper than the $\sim$25 Hz trough (32\% vs.\ 12\% depletion; Table~\ref{tab:aperiodic_null}) and is therefore more reliably detected in individual datasets (Section~\ref{sec:goldenratio}), but the $\sim$25 Hz trough is more stable across smoothing bandwidths. Both are genuine spectral features. The full dataset thus exhibits \emph{hierarchical trough structure}: five boundaries whose inter-trough geometric mean ratio is $\phisym$, with the two shallowest ($\sim$7.6 Hz and $\sim$25 Hz) detectable only at narrow smoothing or large pooled sample sizes. The pooled density is used throughout Part~I to estimate \emph{population geometry}---the positions, depths, and scaling of band boundaries---not to imply that each individual expresses the full five-trough structure with equal clarity. Per-dataset analyses resolve the three coarsest troughs; the two shallower boundaries require pooling for detection, reflecting their dependence on large samples rather than any artefact of pooling. The subject-level bootstrap (Section~\ref{sec:bootstrap}) directly addresses the dependence structure by resampling subjects rather than peaks.

\subsubsection{The Golden Ratio Best Fits Empirical Boundaries}
\label{sec:goldenratio}

Kernel density estimation of all 4,572,489 pooled peaks identifies five density troughs at 5.15, 7.77, 13.43, 25.29, and 35.24 Hz. Consecutive ratios between adjacent troughs are 1.510, 1.728, 1.883, and 1.393. The geometric mean of these four ratios is \textbf{1.6175}, indistinguishable from the golden ratio $\phisym = 1.6180$ to three significant figures.

To determine which candidate scaling ratio best accounts for these boundaries, we fit geometric series models $f_k = \fzero \times r^k$ for seven candidate values of $r$, optimising $\fzero$ by least-squares in log-frequency space (1 free parameter per model), and also tested a free-ratio model (2 free parameters) and linear-Hz spacing (2 parameters). Models were compared by BIC (Table~\ref{tab:model_comparison}).

\begin{table}[h]
\centering
\caption{Model comparison for geometric series fits to empirical density troughs (upper panel) and density peaks (lower panel). BIC is computed with the number of free parameters in parentheses; lower BIC indicates better fit. Mean error is in cents (100 cents = 1 semitone).}
\label{tab:model_comparison}
\begin{threeparttable}
\begin{tabular}{l r r r r r}
\toprule
Model & Ratio $r$ & Optimal $\fzero$ (Hz) & SSE & BIC & Mean error (cents) \\
\midrule
\multicolumn{6}{l}{\textit{Fitting to density troughs (5 boundaries)}} \\
\midrule
$\phisym$ & 1.618 & 5.23 & 0.0257 & $\mathbf{-24.74}$ & \textbf{93} \\
$\sqrt{2}$ & 1.414 & 6.84 & 0.0354 & $-23.14$ & 132 \\
$e - 1$   & 1.718 & 4.63 & 0.0360 & $-23.06$ & 124 \\
$\sqrt{3}$ & 1.732 & 4.56 & 0.0428 & $-22.20$ & 129 \\
$2^{1/3}$ & 1.260 & 9.90 & 0.0557 & $-20.88$ & 140 \\
Octave    & 2.000 & 3.93 & 0.1582 & $-15.66$ & 267 \\
$e$       & 2.718 & 11.20 & 0.2720 & $-12.95$ & 381 \\
Linear-Hz & ---   & ---  & 0.4139 & $-9.24$  & --- \\
Free ratio\tnote{a} & 1.366 & 10.01 & 0.0057 & $-30.69$ & --- \\
\midrule
\multicolumn{6}{l}{\textit{Fitting to density peaks (5 band centres)}} \\
\midrule
$\phisym$ & 1.618 & 6.94 & 0.0287 & $\mathbf{-24.20}$ & 105 \\
$\sqrt{2}$ & 1.414 & 7.41 & 0.0306 & $-23.87$ & 106 \\
$e - 1$   & 1.718 & 3.80 & 0.0571 & $-20.75$ & 164 \\
$2^{1/3}$ & 1.260 & 7.76  & 0.0627 & $-20.29$ & 140 \\
$\sqrt{3}$ & 1.732 & 11.23 & 0.0663 & $-20.00$ & 180 \\
Octave    & 2.000 & 4.25 & 0.0682 & $-19.87$ & 181 \\
$e$       & 2.718 & 9.19 & 0.2026 & $-14.42$ & 327 \\
Linear-Hz & ---   & ---  & 0.4971 & $-8.32$  & --- \\
Free ratio\tnote{a} & 1.334 & 7.50 & 0.0098 & $-27.94$ & --- \\
\bottomrule
\end{tabular}
\begin{tablenotes}
\small
\item[a] Free-ratio model optimises both $\fzero$ and $r$ (2 free parameters).
\end{tablenotes}
\end{threeparttable}
\end{table}

Four findings are clear. First, linear-Hz spacing is excluded with decisive force: $\Delta$BIC $= 15.5$ versus $\phisym$ for troughs, 15.9 for peaks. Second, Euler's number $e$ is strongly excluded (BIC $= -12.95$ vs. $-24.74$ for $\phisym$; $\Delta$BIC $= 11.8$). The \textcite{penttonen2003natural} claim that the inter-band ratio is $\approx e$ is not supported in unbiased peak data. Third, $\phisym$ is the best single fixed-ratio model for both boundary positions (BIC $= -24.74$) and band centres (BIC $= -24.20$), with a mean error of 93 cents (less than one semitone) for troughs. For band centres, $\sqrt{2}$ is close to $\phisym$ (BIC $-23.87$ vs. $-24.20$; $\Delta$BIC $= 0.3$, negligible). Fourth, the discrimination among the three leading candidates---$\phisym$ (1.618), $\sqrt{2}$ (1.414), and $e - 1$ (1.718)---is not sharp: BIC differences of 0.1--1.7 are meaningful but not decisive. All three describe the data far better than any octave-family ratio.

Per-dataset replication yields a different picture. Analysing each of the 9 datasets independently, trough detection reliably finds 2--4 troughs per dataset (typically at $\sim$5.5, $\sim$13.5, and $\sim$35 Hz), with an overall geometric mean ratio of 2.060 (log-SD 0.231). One-sample $t$-tests against candidate constants with $N = 24$ ratios reject $\phisym$ ($t_{23} = +5.02$, $p < 0.001$), $e - 1$ ($p = 0.001$), and $e$ ($p < 0.001$), but cannot reject the octave ratio 2.0 ($t_{23} = +0.61$, $p = 0.548$). This is not a contradiction with the pooled result: individual datasets reliably detect only the three deepest troughs---at $\sim$5, $\sim$13, and $\sim$35 Hz---which are the most depleted boundaries (70\%, 62\%, and 32\% depletion respectively; Table~\ref{tab:aperiodic_null}). These three span ratios of $\sim$2.5, biased toward the octave. The pooled 5-trough solution additionally resolves the two shallower boundaries at $\sim$7.6 Hz (9\% depletion) and $\sim$25 Hz (12\% depletion) that per-dataset analysis misses at typical sample sizes, and those finer boundaries reduce the geometric mean ratio to $\phisym$. The per-dataset octave finding and the pooled $\phisym$ finding are two views of the same hierarchical structure described in Section~\ref{sec:stability}. Crucially, the subject-level bootstrap (Section~\ref{sec:bootstrap}) recovers all five troughs---including the two shallower boundaries at $\sim$7.6 and $\sim$25~Hz---in 100\% of 1,000 iterations, confirming that these finer boundaries are population-level realities detectable with proper dependence-aware resampling, not pooling artifacts that vanish at per-dataset resolution.

\subsubsection{Subject-Level Bootstrap Validates the Five-Trough Solution}
\label{sec:bootstrap}

The preceding analyses pool 4,572,489 peaks and treat them as independent observations. Because peaks are nested within subjects, sessions, and channels, the effective sample size for trough estimation is substantially lower, and confidence intervals from pooled analyses may be artificially narrow. To address this, we performed a stratified bootstrap over subjects: at each of 1,000 iterations, subjects were resampled with replacement independently within each of the 9 datasets (preserving dataset composition), their peaks were pooled, and the full trough-detection pipeline was re-run from scratch.

All five troughs were recovered in 100\% of bootstrap iterations. Position CIs were tight: the $\fzero$ boundary at 7.82~Hz was localized to within $\pm 0.015$~Hz; even the broadest CI (trough 4 at 24.75~Hz: [24.25, 26.01]) spans only 1.75~Hz (Table~\ref{tab:bootstrap_ci}).

\begin{table}[h]
\centering
\caption{Subject-level bootstrap confidence intervals for density trough positions. $N = 2{,}097$ subjects, 1,000 stratified bootstrap iterations. All five troughs detected in 100\% of iterations.}
\label{tab:bootstrap_ci}
\begin{tabular}{r r r r}
\toprule
Trough & Hz & 95\% CI & Width (Hz) \\
\midrule
1 & 5.03  & [4.98, 5.13]   & 0.15 \\
2 & 7.82  & [7.82, 7.85]   & 0.02 \\
3 & 13.59 & [13.40, 13.83] & 0.44 \\
4 & 24.75 & [24.25, 26.01] & 1.75 \\
5 & 34.38 & [34.18, 34.79] & 0.60 \\
\bottomrule
\end{tabular}
\end{table}

The bootstrap distribution of the global inter-trough geometric mean ratio had median 1.617 and 95\% CI [1.609, 1.623] (SD $= 0.004$). Of the candidate constants discussed in the literature (Figure~\ref{fig:bootstrap}):

\begin{itemize}
    \item $\phisym = 1.618$: \textbf{within the CI} ($p = 0.90$)
    \item $e - 1 = 1.718$: excluded ($p < 0.0001$)
    \item $\sqrt{2} = 1.414$: excluded ($p < 0.0001$)
    \item $\sqrt{3} = 1.732$: excluded ($p < 0.0001$)
    \item Octave $= 2.0$: excluded ($p < 0.0001$)
    \item $e = 2.718$: excluded ($p < 0.0001$)
\end{itemize}

The subject-level bootstrap resolves the ambiguity left by the BIC analysis: when the dependence structure is properly accounted for, the 95\% CI on the global scaling ratio spans only 0.014 and contains $\phisym$ while excluding every other named candidate. This does not imply that each adjacent boundary pair is individually $\phisym$-spaced---the local ratios remain heterogeneous (1.36--1.89)---but it establishes that $\phisym$ is the best-supported named global coordinate constant among the theoretical alternatives tested.

\begin{figure}[htbp]
\centering
\includegraphics[width=\textwidth]{fig9_bootstrap.png}
\caption{\textbf{Subject-level bootstrap of trough positions and inter-trough ratio.} (A)~95\% confidence intervals for each of five density trough positions, computed from 1,000 stratified bootstrap resamples over $N = 2{,}097$ subjects (resampled with replacement within each of 9 datasets). Gold dashed lines mark $\phisym$-lattice boundaries. Percentages indicate detection rate (100\% for all troughs). Error bars show 2.5th--97.5th percentile bootstrap CIs. (B)~Bootstrap distribution of the global geometric mean ratio of consecutive trough frequencies (SD $= 0.004$). The 95\% CI [1.609, 1.623] (shaded) contains $\phisym = 1.618$ (gold, two-tailed permutation $p = 0.90$) and excludes all other named candidates (all $p < 0.0001$).}
\label{fig:bootstrap}
\end{figure}

\subsubsection{$\phisym$-Lattice Boundaries are Near-Optimal}
\label{sec:boundaries}

The BIC analysis establishes $\phisym$ as the best fixed-ratio model for empirical boundaries. A separate question is whether the specific boundary \emph{positions} produced by the $\phisym$-lattice ($\fzero \times \phisym^n$ with $\fzero = 7.60$ Hz) are optimal for measurement purposes---that is, whether they produce simple, consistent within-band enrichment profiles. We assess this with a boundary sweep.

A $36 \times 36$ grid of $(\fzero, r)$ pairs spanning $\fzero \in [6.0, 9.5]$ Hz and $r \in [1.3, 2.3]$ was evaluated. For each coordinate system, within-band enrichment profiles were computed and their simplicity measured by the $R^2$ of a low-order polynomial fit. Among all named systems tested (Table~\ref{tab:boundary_sweep}), the $\phisym$-lattice achieves the highest profile simplicity ($R^2 = 0.968$)---enrichment profiles within each band are nearly perfectly described by low-order polynomials when boundaries are placed at $\fzero \times \phisym^n$.

\begin{table}[h]
\centering
\caption{Named boundary systems evaluated in the $36 \times 36$ coordinate sweep.
Profile simplicity is the mean $R^2$ of a low-order polynomial fit to within-band enrichment profiles (higher = simpler, more interpretable profiles). Cross-dataset consistency is the mean pairwise Pearson $r$ of enrichment vectors across 9 datasets.}
\label{tab:boundary_sweep}
\begin{tabular}{l r r r}
\toprule
System & Ratio & Simplicity $R^2$ & Consistency $r$ \\
\midrule
$\phisym$-lattice  & 1.618 & \textbf{0.968} & 0.802 \\
Third-octave       & 1.260 & 0.942          & \textbf{0.896} \\
$\sqrt{2}$         & 1.414 & 0.886          & 0.677 \\
Clinical (4--8--13--30 Hz) & $\sim$1.63 & 0.807 & 0.808 \\
Octave             & 2.000 & 0.647          & 0.819 \\
\bottomrule
\end{tabular}
\end{table}

A boundary slide analysis tested whether each of the four $\phisym$-lattice boundaries is individually near-optimal for profile simplicity (i.e., measurement-optimal, not biologically optimal in any demonstrated sense). Each boundary was independently varied $\pm 25\%$ from its $\phisym$-lattice position while holding the others fixed. All four boundaries lie within 2\% of their simplicity-maximising positions (Table~\ref{tab:boundary_slide}). The theta/alpha boundary at $\fzero = 7.60$ Hz is exactly at the optimum (0.0\% offset).

\begin{table}[h]
\centering
\caption{Boundary slide analysis: optimal boundary positions for profile simplicity relative to $\phisym$-lattice positions.}
\label{tab:boundary_slide}
\begin{tabular}{l r r r}
\toprule
Boundary & $\phisym$-lattice (Hz) & Optimal (Hz) & Offset \\
\midrule
$\theta / \alpha$              & 7.60  & \textbf{7.60}  & \textbf{0.0\%} \\
$\alpha / \beta_{\text{low}}$  & 12.30 & 12.05          & $-2.0\%$ \\
$\beta_{\text{low}} / \beta_{\text{high}}$ & 19.90 & 20.10 & $+1.0\%$ \\
$\beta_{\text{high}} / \gamma$ & 32.19 & 32.52          & $+1.0\%$ \\
\bottomrule
\end{tabular}
\end{table}

\begin{figure}[htbp]
\centering
\includegraphics[width=0.7\textwidth]{fig3_boundary_sweep.png}
\caption{\textbf{Boundary sweep: profile simplicity across coordinate systems.} Heatmap of within-band profile simplicity (mean $R^2$ of low-order polynomial fits to enrichment profiles) across a $36 \times 36$ grid of 1,296 $(\fzero, r)$ coordinate systems, evaluated on 4.57 million peaks from $N = 2{,}097$ subjects. The $\phisym$-lattice (gold star) achieves the highest simplicity ($R^2 = 0.968$). Other named systems are marked for comparison: octave (square), clinical convention (diamond), $\sqrt{2}$ (triangle).}
\label{fig:boundary_sweep}
\end{figure}

\paragraph{Summary.} At the population level, the data support log-frequency organisation unequivocally ($\Delta$BIC $> 46$); support a $\phisym$-like global coordinate constant among named alternatives (bootstrap 95\% CI [1.609, 1.623]); and do \emph{not} support special within-band $\phisym$ landmarks (noble-number positions carry no more signal than simple rationals). The $\phisym$-lattice is a measurement-optimal coordinate system---its value is boundary placement for revealing simple enrichment profiles, not a claim about golden-ratio computation in neural circuits.

\subsubsection{Within Bands, No Fine-Structure at $\phisym$-Specific Positions}
\label{sec:within}

The $\phisym$-lattice places not only boundaries but also interior landmarks within each band: noble numbers, inverse nobles, and the band attractor. These have been proposed as sites of privileged oscillatory activity \parencite{pletzer2010golden}, on the grounds that $\phisym$-irrational positions maximally resist synchronisation with other bands. We test this prediction directly.

Five analyses addressed whether $\phisym$-specific positions carry special signal within bands. First, a landmark capture test compared the enrichment variance explained by $\phisym$-lattice positions against equal-spaced positions and 1000 random position sets: $\phisym$-lattice positions ranked at the 55th--81st percentile of random sets---numerically better than average but not significantly so ($p > 0.05$ in all bands). Second, a feature alignment test asked whether enrichment extrema preferentially fall at $\phisym$ positions: in 3 of 4 testable bands, simple rational positions ($1/2$, $1/3$, $2/3$) were \emph{closer} to the extrema than $\phisym$ positions were. Third, a noble-versus-rational permutation test found no significant enrichment difference between noble numbers and simple rationals in any band ($p = 0.67$--$0.98$). Fourth, periodicity analysis found no sub-octave periodic structure in the within-band enrichment profiles: the dominant period in all bands is 1.0 (the full log-octave), with no harmonics at $1/\phisym$, $1/2$, or other sub-octave values.

Within-band enrichment profiles are smooth, low-dimensional curves---ascending ramps in theta, beta-low, and gamma; a descending ramp in beta-high; a single broad mountain in alpha---fully characterised by 2--3 parameters with no fine-structure at any mathematically special interior positions. The $\phisym$-lattice provides value through its boundary placement, not through any privileged significance of its interior coordinates.

\subsubsection{Relationship to Prior Work}

All results reported in this paper use the v3 extraction pipeline, which corrects three systematic biases identified in a prior report \parencite{lacy2026frontiers}. These corrections---a $\fzero$ alignment mismatch, a Hz-weighting error in the enrichment null hypothesis, and a split-FOOOF boundary artefact---are documented in detail in Section~\ref{sec:methods_corrections} (Methods). The most consequential correction reversed the beta-low boundary enrichment from $+101\%$ to $-33\%$, invalidating the previously reported U-shape profile. Five enrichment values changed sign. The implications for the prior paper's claims are discussed in Section~3.6.
